%----------------Cap_02----------------%

\chapter{Título da Seção Primária}

Esta seção mostra como estruturar um documento em LaTeX.  
O nível mais alto de divisão é a seção primária, criada com o comando:

\begin{verbatim}
\chapter{Seção Primária}
\end{verbatim}

Cada seção primária inicia automaticamente em uma nova página.


\section{Título da Seção Secundária}

Dentro de uma seção primária, podemos criar subdivisões (seções secundárias) usando o seguinte comando:

\begin{verbatim}
\section{Seção Secundária}
\end{verbatim}

Essas seções não iniciam em nova página automaticamente.

Texto texto texto texto texto texto texto texto texto texto texto texto texto texto texto texto texto texto texto texto texto texto texto texto texto texto texto texto texto texto texto texto texto texto texto texto texto texto texto texto texto texto texto texto texto texto texto texto texto texto texto texto texto texto texto texto texto texto texto texto texto texto texto texto texto texto texto texto texto.

\section{Título da Seção Secundária}

Para \cite{ciclano}, ``citação direta, aquela até três linhas, deve aparecer desta forma no texto. Texto texto texto''.

Texto texto texto texto texto texto texto texto texto texto texto texto texto texto texto texto texto texto texto texto texto texto texto texto texto texto texto texto texto texto texto texto texto texto texto texto texto texto texto texto texto texto texto texto texto texto texto texto texto texto texto texto texto texto texto texto texto texto texto texto texto texto texto texto texto texto texto texto texto.

\subsection{Título da seção terciária}

Dentro de uma seção secundária, podemos ainda criar subdivisões menores, chamadas seções terciárias, com o comando:

\begin{verbatim}
\subsection{SEÇÃO TERCIÁRIA}
\end{verbatim}

Elas servem para detalhar melhor a organização do texto em níveis mais específicos.

Texto texto texto texto texto texto texto texto texto texto texto texto texto texto texto texto texto texto texto texto texto texto texto texto texto texto texto texto texto texto texto texto texto texto texto texto texto texto texto texto texto texto texto texto texto texto texto texto texto texto texto texto texto texto texto texto texto texto texto texto texto texto texto texto texto texto texto texto texto.

\subsection{Título da seção terciária}

Texto texto texto texto texto texto texto texto texto texto texto texto texto texto texto texto texto texto texto texto texto texto texto texto texto texto texto texto texto texto texto texto texto texto texto texto texto texto texto texto texto texto texto texto texto texto texto texto texto texto texto texto texto texto texto texto texto texto texto texto texto texto texto texto texto texto texto texto texto.

\subsubsection{Título da seção quaternária}

Texto texto texto texto texto texto texto texto texto texto texto texto texto texto texto texto texto texto texto texto texto texto texto texto texto texto texto texto texto texto texto texto texto texto texto texto texto texto texto texto texto texto texto texto texto texto texto texto texto texto texto texto texto texto texto texto texto texto texto texto texto texto texto texto texto texto texto texto texto.

Texto texto texto texto texto texto texto texto texto texto texto texto texto texto texto texto texto texto texto texto texto texto texto texto texto texto texto texto texto texto texto texto texto texto texto texto texto texto texto texto texto texto texto texto texto texto texto texto texto texto texto texto texto texto texto texto texto texto texto texto texto texto texto texto texto texto texto texto texto.